\subsubsection{Measurement Outcomes and Bob's Corrections}\label{sec:measurement-outcomes-and-bobs-corrections}

Everything so far was just preparation. This is where the teleportation suddenly becomes interesting.

\highspace
Remember where we left off in the previous section: Alice has a composite system of two qubits, the first one being the qubit she wants to teleport and the second one being her half of the shared Bell pair. She has just applied a CNOT gate and a Hadamard gate to her two qubits:
\begin{equation*}
    \begin{aligned}
        \ket{\Psi_{\text{H}}} = \ket{\Psi_{2}} &= \frac{1}{2} \left[ \,
            \ket{00} \otimes \left(\alpha\ket{0} + \beta\ket{1}\right) + \right. \\
            & \hspace{0.88cm} \ket{01} \otimes \left(\beta\ket{0} + \alpha\ket{1}\right) + \\
            & \hspace{0.88cm} \ket{10} \otimes \left(\alpha\ket{0} - \beta\ket{1}\right) + \\
            & \hspace{0.84cm} \left. \ket{11} \otimes \left(-\beta\ket{0} + \alpha\ket{1}\right)
        \, \right]
    \end{aligned}
\end{equation*}
Before Alice measures her two qubits, \textbf{all four branches coexist in superposition}.

\highspace
We have also seen that after measurement, \textbf{only one} of the four branches will \textbf{survive}, and the other three will be destroyed.

\highspace
Let's imagine Alice \hl{measures her two qubits and gets the outcome $\ket{00}$.} In this case, the state immediately becomes:
\begin{equation*}
    \ket{\Psi_3} = \ket{00}\left(\alpha\ket{0} + \beta\ket{1}\right)
\end{equation*}
Now look only at Bob's qubit, which is the last one in the tensor product. It is exactly the \hl{same as the original qubit that Alice wanted to teleport} (state $\ket{\psi} = \alpha\ket{0} + \beta\ket{1}$). In other words, Alice only needs to send $00$ to Bob, and he will know that his qubit is \textbf{already in the correct state} (or formally, he can apply the identity operator to it, which does nothing).

\highspace
\begin{flushleft}
    \textcolor{Green3}{\faIcon{question-circle} \textbf{What if Alice measures a different outcome?}}
\end{flushleft}
Let's complicate things a bit. Suppose Alice measures a different outcome:
\begin{itemize}
    \item If Alice measures $\ket{01}$, the state becomes:
    \begin{equation*}
        \ket{\Psi_3} = \ket{01}\left(\beta\ket{0} + \alpha\ket{1}\right)
    \end{equation*}
    In this case, the amplitudes of Bob's qubit are swapped:
    \begin{equation*}
        \ket{\psi} = \alpha\ket{0} + \beta\ket{1} \longrightarrow \ket{\psi'} = \beta\ket{0} + \alpha\ket{1}
    \end{equation*}
    Therefore, Bob needs to apply an operator that swaps the amplitudes back to their original order. We know that the \hl{Pauli-$X$ gate swaps the amplitudes of a qubit, so Bob can apply it to his qubit to recover the original state:}
    \begin{equation*}
        \ket{\psi'} = \beta\ket{0} + \alpha\ket{1} \xrightarrow{X} \ket{\psi} = \alpha\ket{0} + \beta\ket{1}
    \end{equation*}
    So, if Alice measures $\ket{01}$, she sends $01$ to Bob, and he knows that he needs to apply the Pauli-X gate to his qubit to recover the original state.


    \item If Alice measures $\ket{10}$, the state becomes:
    \begin{equation*}
        \ket{\Psi_3} = \ket{10}\left(\alpha\ket{0} - \beta\ket{1}\right)
    \end{equation*}
    In this case, Bob's qubit has the same amplitudes as the original state, but the amplitude of $\ket{1}$ has a negative sign. Therefore, Bob needs to apply an operator that flips the sign of the amplitude of $\ket{1}$.
    
    For this purpose, we can use the \hl{Pauli-$Z$ gate}, which flips the sign of the amplitude of $\ket{1}$:
    \begin{equation*}
        \ket{\psi'} = \alpha\ket{0} - \beta\ket{1} \xrightarrow{Z} \ket{\psi} = \alpha\ket{0} + \beta\ket{1}
    \end{equation*}
    So, if Alice measures $\ket{10}$, she sends $10$ to Bob, and he knows that he \hl{needs to apply the Pauli-$Z$ gate to his qubit to recover the original state.}


    \item Finally, if Alice measures $\ket{11}$, the state becomes:
    \begin{equation*}
        \ket{\Psi_3} = \ket{11}\left(-\beta\ket{0} + \alpha\ket{1}\right)
    \end{equation*}
    In this case, Bob's qubit has the amplitudes swapped and the amplitude of $\ket{0}$ has a negative sign. Therefore, Bob needs to apply an operator that swaps the amplitudes and flips the sign of the amplitude of $\ket{0}$.

    To fix this issue, Bob can choose two strategies:
    \begin{itemize}
        \item \hl{Apply the Pauli-$X$ gate first to swap the amplitudes, and then apply the Pauli-$Z$ gate to flip the sign of the amplitude of $\ket{0}$} (or vice versa, since the two gates commute):
        \begin{equation*}
            \ket{\psi'} = -\beta\ket{0} + \alpha\ket{1} \xrightarrow{X} \alpha\ket{0} - \beta\ket{1} \xrightarrow{Z} \ket{\psi} = \alpha\ket{0} + \beta\ket{1}
        \end{equation*}

        \item Or simply (cleanly) \hl{apply the Pauli-$Y$ gate}, which swaps the amplitudes and flips the sign of the amplitude of $\ket{0}$ in one step:
        \begin{align*}
            \ket{\psi} &= Y\ket{\psi'} \\
            &= Y(-\beta\ket{0} + \alpha\ket{1}) \\
            &\downarrow \text{ Recall: } Y\ket{0} = i\ket{1}, \quad Y\ket{1} = -i\ket{0} \\
            &= -\beta\left(i\ket{1}\right) + \alpha\left(-i\ket{0}\right) \\
            &= -\beta\left(i\ket{1}\right) - \alpha\left(i\ket{0}\right) \\
            &= -i \underbrace{\left(\alpha\ket{0} + \beta\ket{1}\right)}_{\ket{\psi}} \\
            &\downarrow \text{ Global phase } -i \text{ can be physically ignored} \\
            &= \alpha\ket{0} + \beta\ket{1}
        \end{align*}
    \end{itemize}
\end{itemize}
At the beginning of the protocol we might expect: ``\emph{Bob has absolutely no information about the state of his qubit, so how can he possibly know what to do?}'' But after Alice performs \textbf{only two gates}, Bob's qubit is \textbf{already one of four versions of the correct state}. The quantum information has already reached Bob, but the only thing still missing is \emph{\textbf{which of the four versions is it?}} And that's all the two classical bits communicate. They do \textbf{not} carry the amplitudes $\alpha$ and $\beta$. They only answer the question: ``\emph{Bob, which correction should you apply?}''.